%! Displaying and Analyzing Categorical Data — Unit 3, Lesson 3.1 — Guided Notes & Practice
% THE in-class document, in the MAIN IDEAS / NOTES shape (LESSON_SHAPE.md §1): the vocabulary
% box, the hook, then ONE two-column guidednotes table. Problems run 1..15. Unit 2-on timing
% (5/45/5/5): I do ~25 min (rows 1-4, problems 1-11), Guided Practice ~20 (12-15).
% 4 pages at 12pt — a ceiling (10-page packet, LESSON_SHAPE.md §1).
%   I DO   : the teacher reads the definition, marks up the display, works 1, 3, 6, 8.
%   WE DO  : the class works the rest of each grid; the last row, Guided Practice (12-15), is
%            worked entirely together on the library's streaming survey.
%   YOU DO : nothing on this page — ../homework/ IS the individual practice.
% THE TRAPS: problem 4 (a running total read as a segment) and problem 10 (the wrong whole:
% 54/94 is a percent of bus riders, not of ninth graders).
% THE CRUX: problem 9 (row 4, "Which Whole?"): the hook vote re-taken. More ninth graders ride
% (54 > 40), but a smaller percent of them do (45% < 50%). Problem 15 is the crux transferred.
%
% THE DATA — Oakdale High, 200 students surveyed:
%            Bus  Car  Walk  Total
%   Grade 9   54   48    18    120     45% 40% 15%
%   Grade 12  40   32     8     80     50% 40% 10%
%   Total     94   80    26    200     47% 40% 13%     54/94 = 57.4% of bus riders are ninth graders
% Row 2's segmented bars stack Bus, Car, Walk; running totals 54/102/120 and 40/72/80.
% Row 2 deliberately stays in COUNTS: its segmented bars do not settle the vote. The 100% bars
% appear only in row 4, AFTER problem 9, so the hook stays open until then.
% GUIDED PRACTICE — Fairview Community Library, 250 patrons:
%               StarView ReelBox Pulse Total
%   Under 30        45      30    25   100     45% 30% 25%
%   30 and over     60      66    24   150     40% 44% 16%
%   Total          105      96    49   250     49/250 = 19.6%; 25/49 = 51.0%
%
% PAGE LOCKSTEP with notes_key/: generated by templates/lesson/mkkey.py. KEEP EVERY \pcell ON ONE
% SOURCE LINE. Inside a cell `\\` ENDS THE ROW — break lines with \par. A row cannot break across
% pages: the figure gets its own sub-row and each grid row is its own sub-row (probgrid*).
% NO SKETCH QUESTIONS: every display is pre-drawn in TikZ or set as a table and read.
\documentclass[12pt]{article}
\usepackage{saar-article}
\usepackage{saar-boxes}
\newcommand{\TallMath}[1]{$\displaystyle #1\rule[-1.4em]{0pt}{3.2em}$}
% Word bank strip — ONLY above a table the student fills (row 3). Defined per document;
% the key inherits it and prints the same strip, so heights match.
\newcommand{\wordbank}[1]{\par\vspace{2pt}\noindent\fcolorbox{goldacc}{goldbg}{\parbox{\dimexpr\linewidth-2\fboxsep-2\fboxrule\relax}{\small\textbf{Word bank:} #1}}\par\vspace{4pt}}
% Vocabulary rows: a FIXED-HEIGHT row carrying the term label and OPEN writing space —
% no underline beside the term and no rule line (user correction 2026-09-06). The key fills
% exactly the same height, so the box cannot drift blank vs. keyed. Copied from
% unit01/lesson02/notes/main.tex. Keep key definitions to two lines.
\newcommand{\termrowheight}{2.3cm}
\newlength{\termrowinset}\setlength{\termrowinset}{7pt}
\newcommand{\vterm}[1]{%
  \par\noindent\begin{minipage}[t][\termrowheight][t]{\linewidth}%
    \vspace*{\termrowinset}%
    \textbf{\textcolor{cerulean}{#1:}}%
  \end{minipage}\par}
\newcommand{\vtermans}[2]{%
  \par\noindent\begin{minipage}[t][\termrowheight][t]{\linewidth}%
    \vspace*{\termrowinset}%
    \textbf{\textcolor{cerulean}{#1:}}\quad\ans{#2}%
  \end{minipage}\par}

\begin{document}

\pageheader{Unit 3, Lesson 3.1}{Guided Notes \& Practice}
% NAMESTRIP: no \namedateperiod here. The student writes their name once, on the
% cover this component is stapled behind.

\vspace{0.15in}

\boxguard[16]
\begin{vocabbox}
\small
Fill in each term \textbf{as we name it} in the notes below.
\vterm{Relative frequency}
\vterm{Two-way table}
\vterm{Marginal frequency}
\vterm{Conditional relative frequency}
\end{vocabbox}

\boxguard[12]
\begin{hookbox}
\small
Unit~3 is about data sorted into groups and the chances those data give us. It opens with a
headline. The Oakdale High newsletter reports: \emph{``$54$ ninth graders ride the bus and only
$40$ twelfth graders do, so ninth graders are more likely to ride the bus.''}

\vspace{3pt}
\textbf{Do you agree?} Circle one: \quad Agree \qquad Disagree \qquad Can't tell yet \quad ---
and say why. We vote again at problem~9.
\par\writelines{2}
\end{hookbox}

\small
\begin{guidednotes}
% ── ROW 1: bar graphs of one categorical variable (PS.DS.4a) ─────────────────
\mainidea[One categorical variable]{Bar Graphs} &
A \textbf{bar graph} shows how many individuals fall in each category of one categorical
variable. Its bars can show the \textbf{frequency} (the count) or the \textbf{relative
frequency}: the count divided by the total, as a fraction, decimal, or percent.
\\
& {\centering
{\footnotesize\textbf{How $200$ Oakdale High students get to school (spreadsheet charts)}}\par\vspace{4pt}
\begin{tikzpicture}[x=1cm, y=0.03cm]
  % frequency panel
  \foreach \v in {20,40,...,100} { \draw[linegray] (0,\v) -- (4.6,\v); }
  \foreach \v in {0,20,...,100} { \node[left, font=\scriptsize] at (0,\v) {\v}; }
  \draw[thick] (0,0) -- (0,104);
  \draw[thick] (0,0) -- (4.6,0);
  \fill[cerulean, draw=charcoal] (0.45,0) rectangle (1.25,94);
  \fill[frostmid, draw=charcoal] (1.9,0) rectangle (2.7,80);
  \fill[goldbg, draw=charcoal] (3.35,0) rectangle (4.15,26);
  \node[above, font=\scriptsize\bfseries] at (0.85,94) {94};
  \node[above, font=\scriptsize\bfseries] at (2.3,80) {80};
  \node[above, font=\scriptsize\bfseries] at (3.75,26) {26};
  \node[below, font=\scriptsize] at (0.85,0) {Bus};
  \node[below, font=\scriptsize] at (2.3,0) {Car};
  \node[below, font=\scriptsize] at (3.75,0) {Walk};
  \node[rotate=90, font=\scriptsize\itshape] at (-0.85,52) {number of students};
  % relative frequency panel
  \begin{scope}[xshift=6.7cm]
  \foreach \v in {10,20,...,50} { \draw[linegray] (0,{\v*2}) -- (4.6,{\v*2}); }
  \foreach \v in {0,10,...,50} { \node[left, font=\scriptsize] at (0,{\v*2}) {\v\%}; }
  \draw[thick] (0,0) -- (0,104);
  \draw[thick] (0,0) -- (4.6,0);
  \fill[cerulean, draw=charcoal] (0.45,0) rectangle (1.25,94);
  \fill[frostmid, draw=charcoal] (1.9,0) rectangle (2.7,80);
  \fill[goldbg, draw=charcoal] (3.35,0) rectangle (4.15,26);
  \node[below, font=\scriptsize] at (0.85,0) {Bus};
  \node[below, font=\scriptsize] at (2.3,0) {Car};
  \node[below, font=\scriptsize] at (3.75,0) {Walk};
  \node[rotate=90, font=\scriptsize\itshape] at (-0.95,52) {percent of students};
  \end{scope}
\end{tikzpicture}\par\vspace{5pt}
Left: \labelbox{2.0cm}{} bar graph\quad Right: \labelbox{3.5cm}{} bar graph\par}
\\
& \notesprompt{Read each graph --- every answer names the students and the category.}
\begin{probgrid}{|Y|Y|}
\pcell{1}{Read the left graph. How many of the surveyed students ride the bus?}{2.0cm}{} & \pcell{2}{Find the relative frequency of Car as a fraction and a percent. Check it on the right graph, then say what it means.}{2.0cm}{} \\ \hline
\end{probgrid}
\\ \hline
% ── ROW 2: side-by-side and segmented bars (PS.DS.4b) ────────────────────────
\mainidea[Comparing groups]{Two Kinds of Bars} &
\textbf{Side-by-side bars} stand the groups next to each other within each category.
\textbf{Segmented bars} stack each group's categories into one bar, so you see the group's
whole and its parts. Both graphs here show counts.
\\
& {\centering
\begin{tikzpicture}[x=1cm, y=0.05cm]
  % side-by-side panel (counts)
  \foreach \v in {20,40,60} { \draw[linegray] (0,\v) -- (5.2,\v); }
  \foreach \v in {0,20,40,60} { \node[left, font=\scriptsize] at (0,\v) {\v}; }
  \draw[thick] (0,0) -- (0,63);
  \draw[thick] (0,0) -- (5.2,0);
  \foreach \c/\a/\b in {0.3/54/40, 2.0/48/32, 3.7/18/8} {
    \fill[cerulean, draw=charcoal] (\c,0) rectangle ({\c+0.6},\a);
    \fill[goldbg, draw=charcoal] ({\c+0.6},0) rectangle ({\c+1.2},\b);
    \node[above, font=\scriptsize\bfseries] at ({\c+0.3},\a) {\a};
    \node[above, font=\scriptsize\bfseries] at ({\c+0.9},\b) {\b};
  }
  \node[below, font=\scriptsize] at (0.9,0) {Bus};
  \node[below, font=\scriptsize] at (2.6,0) {Car};
  \node[below, font=\scriptsize] at (4.3,0) {Walk};
  \fill[cerulean, draw=charcoal] (0.5,70) rectangle (0.75,75);
  \node[right, font=\scriptsize] at (0.75,72.5) {Grade 9};
  \fill[goldbg, draw=charcoal] (2.6,70) rectangle (2.85,75);
  \node[right, font=\scriptsize] at (2.85,72.5) {Grade 12};
  \node[font=\scriptsize\bfseries] at (2.6,84) {Side-by-side};
  % segmented panel (counts), running totals beside each bar
  \begin{scope}[xshift=6.9cm, y=0.0275cm]
  \foreach \v in {40,80,120} { \draw[linegray] (0,\v) -- (4.8,\v); }
  \foreach \v in {0,40,80,120} { \node[left, font=\scriptsize] at (0,\v) {\v}; }
  \draw[thick] (0,0) -- (0,122);
  \draw[thick] (0,0) -- (4.8,0);
  \fill[cerulean, draw=charcoal] (0.5,0) rectangle (1.5,54);
  \fill[frostmid, draw=charcoal] (0.5,54) rectangle (1.5,102);
  \fill[goldbg, draw=charcoal] (0.5,102) rectangle (1.5,120);
  \fill[cerulean, draw=charcoal] (2.9,0) rectangle (3.9,40);
  \fill[frostmid, draw=charcoal] (2.9,40) rectangle (3.9,72);
  \fill[goldbg, draw=charcoal] (2.9,72) rectangle (3.9,80);
  \foreach \x/\t in {1.5/54, 1.5/102, 3.9/40, 3.9/72} {
    \node[right, font=\scriptsize\bfseries, inner sep=1.5pt] at (\x,\t) {\t}; }
  \node[above, font=\scriptsize\bfseries, inner sep=1.5pt] at (1.0,120) {120};
  \node[above, font=\scriptsize\bfseries, inner sep=1.5pt] at (3.4,80) {80};
  \node[below, font=\scriptsize] at (1.0,0) {Grade 9};
  \node[below, font=\scriptsize] at (3.4,0) {Grade 12};
  \fill[cerulean, draw=charcoal] (0.2,138) rectangle (0.5,147);
  \node[right, font=\scriptsize] at (0.5,142.5) {Bus};
  \fill[frostmid, draw=charcoal] (1.6,138) rectangle (1.9,147);
  \node[right, font=\scriptsize] at (1.9,142.5) {Car};
  \fill[goldbg, draw=charcoal] (3.0,138) rectangle (3.3,147);
  \node[right, font=\scriptsize] at (3.3,142.5) {Walk};
  \node[font=\scriptsize\bfseries] at (2.4,160) {Segmented (running totals)};
  \end{scope}
\end{tikzpicture}\par}
\\
& \notesprompt{Read the bars --- every answer names the grade and the count.}
\begin{probgrid}{|Y|Y|}
\pcell{3}{Side-by-side: in which category is the gap between the two grades' counts largest? How large is it?}{2.0cm}{} & \pcell{4}{Segmented: how many twelfth graders come by car? Why is the answer not $72$?}{2.0cm}{} \\ \hline
\end{probgrid}
\\ \hline
% ── ROW 3: generate a two-way table from a summary (PS.DS.4c) ────────────────
\mainidea[From a summary]{Two-Way Tables} &
A \textbf{two-way table} records two categorical variables at once, one in the rows and one in
the columns. Each inside cell is a \textbf{joint} count; each Total is a \textbf{marginal
frequency}.
\par\vspace{4pt}
\fcolorbox{cerulean}{frost}{\parbox{\dimexpr\linewidth-2\fboxsep-2\fboxrule\relax}{\raggedright
\textbf{The summary:} Oakdale surveyed $120$ ninth graders and $80$ twelfth graders. In Grade~$9$,
$54$ ride the bus and $48$ come by car. In Grade~$12$, $40$ ride the bus and $32$ come by car.
Everyone else walks.}}
\\
& \textbf{5.} Build the table from the summary. Subtract to find what the summary leaves out.
\wordbank{8 \;\textbullet\; 18 \;\textbullet\; 26 \;\textbullet\; 32 \;\textbullet\; 40 \;\textbullet\; 48 \;\textbullet\; 54 \;\textbullet\; 80 \;\textbullet\; 94 \;\textbullet\; 120 \;\textbullet\; 200}
{\renewcommand{\arraystretch}{1.5}%
\begin{tabularx}{\linewidth}{@{} l Y Y Y Y @{}}
\toprule
 & \textbf{Bus} & \textbf{Car} & \textbf{Walk} & \textbf{Total} \\
\midrule
\textbf{Grade 9} & \blank{1.2cm} & \blank{1.2cm} & \blank{1.2cm} & \blank{1.2cm} \\
\textbf{Grade 12} & \blank{1.2cm} & \blank{1.2cm} & \blank{1.2cm} & \blank{1.2cm} \\
\textbf{Total} & \blank{1.2cm} & \blank{1.2cm} & \blank{1.2cm} & \blank{1.2cm} \\
\bottomrule
\end{tabularx}}
\\
& \notesprompt{Read the table --- name the kind of number you used.}
\begin{probgrid}{|Y|Y|}
\pcell{6}{How many twelfth graders walk to school? Joint count or marginal frequency?}{2.0cm}{} & \pcell{7}{How many surveyed students come by car? Joint count or marginal frequency?}{2.0cm}{} \\ \hline
\end{probgrid}
\\ \hline
% ── ROW 4: THE CRUX — marginal, relative, conditional (PS.DS.4d) ─────────────
\mainidea[Marginal, relative, conditional]{Which Whole?} &
A \textbf{marginal relative frequency} divides a Total by the whole table. A
\textbf{conditional relative frequency} divides an inside cell by the total of \emph{one} group.
\textbf{In ``Of the [group], what percent \ldots?'' the words after \emph{of} name the whole, and
that total is the denominator.}
\\
& \begin{probgrid}{|Y|Y|}
\pcell{8}{What percent of all $200$ surveyed students ride the bus? Name the kind of frequency.}{2.2cm}{} & \pcell{9}{\textbf{Vote again:} Agree \; Disagree. Find the percent of ninth graders who ride the bus and the percent of twelfth graders who do. Who is more likely?}{2.2cm}{} \\ \hline
\end{probgrid}
\\
& {\centering
{\footnotesize\textbf{The same survey, each grade drawn as its own $100\%$}}\par\vspace{4pt}
\begin{tikzpicture}[x=0.1cm, y=1cm]
  \foreach \v in {0,20,...,100} {
    \draw[linegray] (\v,0) -- (\v,1.75);
    \node[below, font=\scriptsize] at (\v,0) {\v\%}; }
  \draw[thick] (0,0) -- (100,0);
  \foreach \y/\g/\a/\b/\c in {1.05/Grade 9/45/40/15, 0.2/Grade 12/50/40/10} {
    \fill[cerulean, draw=charcoal] (0,\y) rectangle (\a,\y+0.55);
    \fill[frostmid, draw=charcoal] (\a,\y) rectangle ({\a+\b},\y+0.55);
    \fill[goldbg, draw=charcoal] ({\a+\b},\y) rectangle (100,\y+0.55);
    \node[font=\footnotesize\bfseries, text=white] at ({\a/2},\y+0.275) {Bus \a\%};
    \node[font=\footnotesize\bfseries] at ({\a+\b/2},\y+0.275) {Car \b\%};
    \node[font=\tiny\bfseries, inner sep=0pt] at ({\a+\b+\c/2},\y+0.39) {Walk};
    \node[font=\tiny\bfseries, inner sep=0pt] at ({\a+\b+\c/2},\y+0.16) {\c\%};
    \node[anchor=east, font=\scriptsize\bfseries] at (-1,\y+0.275) {\g};
  }
\end{tikzpicture}\par\vspace{5pt}
\fcolorbox{cerulean}{frost}{\parbox{\dimexpr\linewidth-2\fboxsep-2\fboxrule\relax}{\centering
\textbf{Counts say how many. Conditional percents say how likely. To compare groups of different
sizes, compare conditional percents.}}}\par}
\\
& \begin{probgrid}{|Y|Y|}
\pcell{10}{Dani says, \emph{``$54/94 \approx 57\%$, so $57\%$ of ninth graders ride the bus.''} What whole did she use? What does $57\%$ really tell you?}{2.2cm}{} & \pcell{11}{More ninth graders than twelfth graders come by car ($48$ to $32$). Are ninth graders more likely to come by car? Show the percents.}{2.2cm}{} \\ \hline
\end{probgrid}
\\ \hline
% ── ROW 5: GUIDED PRACTICE — we do, together, a fresh context ────────────────
\mainidea[Guided practice]{Streaming Choices} &
The Fairview Community Library asked $250$ patrons which streaming service they prefer. Here is
the library's spreadsheet.
\par\vspace{4pt}
{\centering
{\renewcommand{\arraystretch}{1.3}\ttfamily
\begin{tabular}{|>{\columncolor{graybg}}c|l|c|c|c|c|}
\hline
\rowcolor{graybg} & A & B & C & D & E \\ \hline
1 & \textbf{Age group} & \textbf{StarView} & \textbf{ReelBox} & \textbf{Pulse} & \textbf{Total} \\ \hline
2 & Under 30 & 45 & 30 & 25 & 100 \\ \hline
3 & 30 and over & 60 & 66 & 24 & 150 \\ \hline
4 & \textbf{Total} & 105 & 96 & 49 & 250 \\ \hline
\end{tabular}}\par}
\\
& \rule{0pt}{16pt}\textbf{We work these together.}\par\vspace{3pt}
\begin{probgrid}{|Y|Y|}
\pcell{12}{What percent of all $250$ patrons prefer Pulse? Name the kind of frequency.}{2.2cm}{} & \pcell{13}{Of the patrons under $30$, what percent prefer Pulse? Name the kind of frequency.}{2.2cm}{} \\ \hline
\end{probgrid}
\\
& \begin{probgrid}*{|Y|Y|}
\pcell{14}{Ty says, \emph{``$25/49 \approx 51\%$, so $51\%$ of patrons under $30$ prefer Pulse.''} What whole did he use? Write the true sentence.}{2.4cm}{} & \pcell{15}{The flyer says, \emph{``More patrons $30$ and over chose StarView ($60$ to $45$), so they are more likely to prefer it.''} Check it with percents.}{2.4cm}{} \\ \hline
\end{probgrid}
\\ \hline
\end{guidednotes}

% ── THE NOTES END AT GUIDED PRACTICE ─────────────────────────────────────────
% There is deliberately no "you do" here: ../homework/ is the individual practice,
% assigned at Close & Assign and done outside class.

\end{document}
